\BourbakiExerciseGroup

\begin{enumerate}
\def\labelenumi{\arabic{enumi})}
\tightlist
\item
  The polynomial \(X^{2}-2\) is irreducible in Q{[}X{]} (cf.~I, p.~51, Th. 7); let \(\alpha\) be one of its roots in an algebraic closure \(\Omega\) of Q. Show that the polynomial \(X^{2}-\alpha\) is irreducible over \(E = Q(\alpha)\) ; let \(\beta\) be one of the roots of this polynomial in \(\Omega\) and let \(F = E(\beta) = Q(\beta)\) . Show that F is not a quasi-Galois extension of Q; what is the quasi-Galois extension of Q generated by F?
\end{enumerate}

* 2) Let E be the relative algebraic closure of Q in R (« the field of real algebraic numbers »); if u is an automorphism of E, then \(u(x) = x\) for all \(x \in \mathbf{Q}\) ; show that if \(y \in E\) is such that \(y > 0\) , then also \(u(y) > 0\) . Deduce that \(u(y) = y\) for all \(y \in E\) . *

\begin{enumerate}
\def\labelenumi{\arabic{enumi})}
\setcounter{enumi}{2}
\item
  Show that every algebraic extension of a field \(\mathbf{K}\) , generated by a set of elements of degree 2 over \(\mathbf{K}\) is quasi-Galois over \(\mathbf{K}\) .
\item
  Let K be a field, f an irreducible polynomial in K{[}X{]}, separable over K and of degree n and let \(\alpha_{i}\) ( \(1 \leqslant i \leqslant n\) ) be its roots in an algebraic closure \(\Omega\) of K. Let g be a polynomial in K{[}X{]}, and h an irreducible factor in K{[}X{]} of the polynomial \(f(g(\mathbf{X}))\) . Show that the degree of h is a multiple rn of n and that h has exactly r roots in common with each of the polynomials \(g(\mathbf{X}) - \alpha_{i}\) ( \(1 \leqslant i \leqslant n\) ) of \(\Omega[\mathbf{X}]\) (consider the conjugates over K of a root of h).
\item
  \begin{enumerate}
  \def\labelenumii{\alph{enumii})}
  \tightlist
  \item
    Let K be a field, \(\Omega\) an algebraic closure of K and E, F two subextensions of \(\Omega\) . Show that every composite extension of E and F (V, p.~12) is isomorphic to a composite extension of the form \((L_{w}, 1, w)\) , where w denotes a K-isomorphism of F onto a subfield of \(\Omega\) , 1 is the identity mapping of E and \(L_{w}\) is the subfield of \(\Omega\) generated by \(E \cup w(F)\) .
  \end{enumerate}
\end{enumerate}

\begin{enumerate}
\def\labelenumi{\alph{enumi})}
\setcounter{enumi}{1}
\item
  Deduce from a) that if F is of finite degree n over K, there are at most n composite extensions of E and F that are pairwise non-isomorphic.
\item
  Suppose that E is a quasi-Galois extension of K and F a subextension of E. Show that every composite extension of E and F is isomorphic to a composite extension of the form \((E, 1, v)\) , where v is a K-isomorphism of F onto a subfield of E.
\end{enumerate}

\begin{enumerate}
\def\labelenumi{\arabic{enumi})}
\setcounter{enumi}{5}
\tightlist
\item
  \begin{enumerate}
  \def\labelenumii{\alph{enumii})}
  \tightlist
  \item
    Let K be a field of characteristic exponent p, L an extension of K, \(\Omega\) an algebraic closure of L and \(\bar{K}\) the relative algebraic closure of K in \(\Omega\) . Show that \(\bar{K}\) and \(L^{p^{-\infty}}\) are linearly disjoint over \(E = \bar{K} \cap L^{p^{-\infty}}\) . (Let \((a_{\lambda})\) be a basis of the vector space \(L^{p^{-\infty}}\) over E; consider a relation \(\sum_{\lambda} c_{\lambda} a_{\lambda} = 0\) , where \(c_{\lambda} \in \bar{K}\) and the number of non-zero
  \end{enumerate}
\end{enumerate}

\(c_{\lambda}\) is as small as possible, and apply an L-automorphism of \(\Omega\) to this relation.)

\begin{enumerate}
\def\labelenumi{\alph{enumi})}
\setcounter{enumi}{1}
\tightlist
\item
  Deduce from a) that if E and F are two extensions of K such that E is algebraic over K and K is relatively algebraically closed in F, then any two composite extensions of E and F are isomorphic.
\end{enumerate}

\begin{enumerate}
\def\labelenumi{\arabic{enumi})}
\setcounter{enumi}{6}
\tightlist
\item
  Let K be a field of characteristic 0, \(\alpha\) , \(\beta\) , \(\xi\) three elements of an algebraic closure of K such that \([\mathbf{K}(\alpha):\mathbf{K}] = 2\) , \([\mathbf{K}(\beta):\mathbf{K}] = 3\) , \([\mathbf{K}(\xi):\mathbf{K}] = n > 1\) ; denote by \(\alpha'\) (resp. \(\beta'\) , \(\beta''\) ) the unique conjugate \(\neq \alpha\) of \(\alpha\) (resp. the conjugates \(\neq \beta\) of \(\beta\) ).
\end{enumerate}

\begin{enumerate}
\def\labelenumi{\alph{enumi})}
\item
  Show that \([\mathbf{K}(\alpha, \xi): \mathbf{K}] = n\) if \(\alpha \in \mathbf{K}(\xi)\) , \([\mathbf{K}(\alpha, \xi): \mathbf{K}] = 2n\) otherwise.
\item
  Show that \([\mathbf{K}(\beta,\xi):\mathbf{K}]=n\) if \(\beta\in\mathbf{K}(\xi)\) , \([\mathbf{K}(\beta,\xi):\mathbf{K}]=2n\) if \(\beta\in\mathbf{K}(\xi)\) and just one of \(\beta'\) , \(\beta''\) belongs to \(\mathbf{K}(\xi)\) and finally \([\mathbf{K}(\beta,\xi):\mathbf{K}]=3n\) if none of the three elements \(\beta\) , \(\beta'\) , \(\beta''\) belongs to \(\mathbf{K}(\xi)\) .
\item
  If \(d\) is the discriminant of the minimal polynomial of \(\beta\) , show that \(\mathbf{K}(\beta, \beta') = \mathbf{K}(\beta, \sqrt{d})\) .
\end{enumerate}

\(d\) ) Show that \(\mathbf{K}(\beta, \beta') = \mathbf{K}(\beta - \beta')\) .

\begin{enumerate}
\def\labelenumi{\alph{enumi})}
\setcounter{enumi}{4}
\item
  If \(n = 2\) and \(\alpha \notin \mathbf{K}(\xi)\) , show that \(\mathbf{K}(\alpha, \xi) = \mathbf{K}(\alpha + \xi)\) .
\item
  Show that \(\mathbf{K}(\alpha, \beta) = \mathbf{K}(\alpha + \beta)\) .
\end{enumerate}

\(g\) ) Show that \(\mathbf{K}(\alpha ,\beta) = \mathbf{K}(\alpha \beta)\) except when \(\alpha = aj\) and \(\beta^3 = b\) , where \(a,b\) are elements of \(\mathbf{K}\) and \(j^3 = 1\) .

\begin{enumerate}
\def\labelenumi{\arabic{enumi})}
\setcounter{enumi}{7}
\tightlist
\item
  Let K be an infinite field, E an algebraic extension of K of finite degree and \(\tau\) an automorphism of the vector K-space E such that for all \(x \in E\) , \(\tau(x)\) is conjugate to x over K. Show that \(\tau\) is a K-automorphism of the field E. (Let N be the quasi-Galois extension of K generated by E and let \(\Gamma\) be the group of K-automorphisms of N, \(\Delta\) the subgroup of \(\Gamma\) consisting of all E-automorphisms and \(\sigma_{j}\) ( \(1 \leqslant j \leqslant r\) ) a system of representatives of the left cosets of \(\Delta\) . If \((u_{i})_{1 \leqslant i \leqslant n}\) is a basis of the vector K-space E, consider the polynomial
\end{enumerate}

\[
\mathrm{F} \left(\mathrm{X} _ {1}, \dots , \mathrm{X} _ {n}\right) = \prod_ {j = 1} ^ {r} \left(\sum_ {i = 1} ^ {n} \left(\tau \left(u _ {i}\right) - \sigma_ {j} \left(u _ {i}\right)\right) \mathrm{X} _ {i}\right).
\]

